% Custo de saída contra validade — o argumento econômico da tese em um plano.
% Rótulos posicionados um a um: `nodes near coords` colidia nos aglomerados
% (A1/A1g à direita, A2/A2g/A3 à esquerda). Eixo x logarítmico porque os
% tokens variam de 138 a 2.041.
\begin{figure}[htb]
	\Caption{\label{fig:custo} Custo de saída contra validade sintática, por braço}
	\centering
	\begin{tikzpicture}
		\begin{axis}[
			width=0.86\textwidth, height=6.8cm,
			xmode=log, log basis x=10,
			xlabel={\textit{Tokens} emitidos por item (mediana, escala log)},
			ylabel={Saídas válidas no XSD (\%)},
			xmin=105, xmax=3200, ymin=-8, ymax=115,
			xtick={125,250,500,1000,2000}, xticklabels={125,250,500,1000,2000},
			ytick={0,25,50,75,100},
			ymajorgrids, xmajorgrids, grid style={draw=black!12},
			axis lines*=left, tick align=outside,
			tick label style={font=\small}, label style={font=\small},
			legend style={font=\small, at={(0.97,0.52)}, anchor=east, draw=none, fill=white, fill opacity=0.9, text opacity=1},
			legend cell align=left,
		]
			\addplot[only marks, mark=*, mark size=2.8pt, corXML]
				coordinates {(631,98.1) (699,94.3)};
			\addlegendentry{XML direto}
			\addplot[only marks, mark=square*, mark size=2.8pt, corDSL]
				coordinates {(138,64.8) (149,54.7) (159,58.5) (2041,0.0)};
			\addlegendentry{DSL por \textit{prompting}}
			\addplot[only marks, mark=triangle*, mark size=4pt, corSFT]
				coordinates {(200,86.8)};
			\addlegendentry{DSL, 7B com ajuste fino}

			% Rótulos individuais, ancorados para evitar sobreposição
			\node[font=\scriptsize, anchor=south]      at (axis cs:631,99.5)  {A1};
			\node[font=\scriptsize, anchor=north west] at (axis cs:699,92.5)  {A1g};
			\node[font=\scriptsize, anchor=south east] at (axis cs:200,88.5)  {\textbf{A4}};
			\node[font=\scriptsize, anchor=south]      at (axis cs:138,66.5)  {A2};
			\node[font=\scriptsize, anchor=west]       at (axis cs:159,58.5)  {\ A3};
			\node[font=\scriptsize, anchor=north]      at (axis cs:149,52.5)  {A2g};
			\node[font=\scriptsize, anchor=south east] at (axis cs:2041,2.5)  {A3m};
		\end{axis}
	\end{tikzpicture}
	\Fonte{Elaborado pelo autor}
	\Nota{Canto superior esquerdo é o desejável: alta validade a baixo custo. O
	braço A4 emite 200 \textit{tokens} medianos contra 631 do A1, mantendo 86{,}8\%
	de validade, e executa em placa de consumidor. O A3m aparece isolado à direita
	por gerar XML sob teto dimensionado para DSL, truncando em metade das
	execuções.}
\end{figure}
